\subsection{Demostración aislada por sensor y tipo de dispositivo}
La identidad física de los productores de Madrid v2 no está acreditada. Para
demostrar MG01 y MG07 sin atribuirles dispositivos ficticios, se utiliza una base
académica independiente con versión \archivo{smartcity-academic-v1}. Contiene dos
tipos de dispositivo, cuatro sensores, nueve lecturas y tres documentos narrativos,
todos sintéticos. Las identidades de los sensores proceden del paquete
administrativo de demostración. Los tipos reproducen la definición administrativa
del modelo relacional, sin consultar ni modificar PostgreSQL \cite{laboratorio}.

\begin{figure}[ht]
\centering
\begin{tikzpicture}[box/.style={draw,rounded corners,align=center,font=\small,
  text width=3.1cm,minimum height=1.1cm},ref/.style={-{Latex[length=2mm]}}]
\node[box] (r) {readings\\sensor\_id; valor; magnitud; unidad; fecha BSON};
\node[box,right=1.2cm of r] (s) {sensors\\sensor\_id; sensor\_type; procedencia};
\node[box,right=1.2cm of s] (t) {sensor\_types\\identificador; nombre; categoría};
\node[box,text width=3.7cm,below=1.3cm of s] (n) {dynamic\_information\\texto libre; atributos variables; revisión; fecha BSON};
\draw[ref] (r) -- (s);
\draw[ref] (s) -- (t);
\draw[ref] (n) -- node[right,font=\scriptsize]{sujeto ficticio} (s);
\end{tikzpicture}
\caption{Laboratorio sintético separado de las siete colecciones Madrid v2.
Las referencias se comprueban en la aplicación; no son FK MongoDB.}
\end{figure}

MG01 recibe un \archivo{sensor_id} explícito y recupera sus documentos completos,
ordenados por fecha e identificador. El sensor A1 tiene dos valores de prueba:
10 y 14 grados Celsius. Un identificador inexistente devuelve cero documentos;
la omisión del parámetro se rechaza. No se renombra un canal como sensor.

MG07 usa dos \texttt{\$lookup} para relacionar lecturas, sensores y tipos. El
promedio se pondera por lectura y se calcula por tipo, magnitud y unidad, sin
promediar unidades incompatibles ni promedios de sensores. A y B son alias de dos
tipos administrativos sintéticos, no magnitudes ni modelos físicos observados.

\begin{center}
\begin{tabular}{llllr}
\toprule Tipo & Magnitud & Unidad & Lecturas & Promedio \\
\midrule
A & Temperatura & °C & 4 & 17 \\
B & Temperatura & °C & 3 & $100/3$ \\
B & Humedad relativa & \% & 2 & 50 \\
\bottomrule
\end{tabular}
\end{center}

Las funciones ejecutables están en
\archivo{sprint1/mongodb/academico/consultas_demo.js}. Las consultas públicas
anteriores por canal y por fuente/magnitud/unidad permanecen como funcionalidades
adicionales con sus evidencias originales. Las nueve lecturas de prueba no se
suman a las 748.889 observaciones públicas.

\subsection{Información dinámica y texto no estructurado}
Tres documentos narran escenarios ficticios de inspección y percepción ciudadana.
El texto libre conserva el contexto descriptivo sin reducirlo a columnas de
medición. El sobre tiene identificador, sujeto ficticio, fecha BSON, revisión y
origen sintético; el objeto de atributos admite estructuras heterogéneas.

La búsqueda del término ``vibración'' recupera dos documentos. La consulta por
existencia del atributo de mantenimiento devuelve uno. La selección por sujeto
recupera las revisiones 1 y 2 de A1, conservadas como documentos independientes.
Un tercer relato utiliza atributos diferentes. El contenido narrativo es no
estructurado; el sobre y sus atributos son estructurados o semiestructurados.
No se equipara almacenamiento BSON con información no estructurada.

La ejecución nativa y el contraste Python verifican los documentos completos,
los valores agregados y el contenido de las selecciones. La demostración acredita
almacenamiento y recuperación de textos sintéticos; no incorpora reportes
ciudadanos reales ni observaciones municipales nuevas.
